% Série homocedástica (simulada) × heterocedástica (PETR4), Cap. I.
% Dados de dados/resultados/heterocedasticidade_2026-11-02/serie.csv
% (ferramentas/scripts/heterocedasticidade.sh). As duas partes na mesma escala.
% Compilar: pdflatex heterocedasticidade.tex
\documentclass[border=4pt]{standalone}
\usepackage[utf8]{inputenc}
\usepackage[T1]{fontenc}
\usepackage{mathptmx}
\usepackage{pgfplots}
\usepgfplotslibrary{dateplot,groupplots}
\pgfplotsset{compat=1.18}
\definecolor{csim}{RGB}{120,120,120}
\definecolor{cpetr}{RGB}{0,90,181}
\newcommand{\arq}{../../dados/resultados/heterocedasticidade_2026-11-02/serie.csv}
\begin{document}
\begin{tikzpicture}
\begin{groupplot}[
  group style={group size=1 by 2, vertical sep=0.9cm, xticklabels at=edge bottom},
  width=14.6cm, height=4.6cm, date coordinates in=x,
  xticklabel={\month/\year}, max space between ticks=80,
  xticklabel style={font=\normalsize}, yticklabel style={font=\normalsize},
  xmin=2019-01-02, xmax=2021-12-30, ymin=-37, ymax=22, ytick={-30,-20,-10,0,10,20},
  ylabel={Retorno (\%)}, table/col sep=comma,
  grid=major, grid style={gray!20}, every axis plot/.append style={line width=0.5pt},
  title style={font=\normalsize, yshift=-4pt},
]
\nextgroupplot[title={Série simulada com variância constante (homocedástica)}]
\addplot[csim] table[x=data, y=simulada]{\arq};
\nextgroupplot[title={Retorno diário da PETR4 (heterocedástica)}]
\addplot[cpetr] table[x=data, y=petr4]{\arq};
\end{groupplot}
\end{tikzpicture}
\end{document}
